\documentclass{../base/canon}

%% canon_variables.tex — Valores por defecto de metadatos institucionales
%%
%% Incluir ANTES del \begin{document} para sobreescribir los defaults de la clase.
%% El template puede sobreescribir cualquier valor después de este \input.
%%
%% Uso en template:
%%   %% canon_variables.tex — Valores por defecto de metadatos institucionales
%%
%% Incluir ANTES del \begin{document} para sobreescribir los defaults de la clase.
%% El template puede sobreescribir cualquier valor después de este \input.
%%
%% Uso en template:
%%   %% canon_variables.tex — Valores por defecto de metadatos institucionales
%%
%% Incluir ANTES del \begin{document} para sobreescribir los defaults de la clase.
%% El template puede sobreescribir cualquier valor después de este \input.
%%
%% Uso en template:
%%   \input{../base/canon_variables}
%%   \SetDocTitle{Mi Informe Específico}   % sobreescribe el default

\SetDocLogo{../assets/logo.pdf}
\SetDocUnit{Tecnología \& Sistemas}
\SetContactUnit{Mesa de soporte}
\SetContactMembers{%
  Equipo CoreX & Soporte & soporte@empresa.com%
}
\SetContactNote{Para soporte técnico o consultas sobre este documento.}

%%   \SetDocTitle{Mi Informe Específico}   % sobreescribe el default

\SetDocLogo{../assets/logo.pdf}
\SetDocUnit{Tecnología \& Sistemas}
\SetContactUnit{Mesa de soporte}
\SetContactMembers{%
  Equipo CoreX & Soporte & soporte@empresa.com%
}
\SetContactNote{Para soporte técnico o consultas sobre este documento.}

%%   \SetDocTitle{Mi Informe Específico}   % sobreescribe el default

\SetDocLogo{../assets/logo.pdf}
\SetDocUnit{Tecnología \& Sistemas}
\SetContactUnit{Mesa de soporte}
\SetContactMembers{%
  Equipo CoreX & Soporte & soporte@empresa.com%
}
\SetContactNote{Para soporte técnico o consultas sobre este documento.}

\SetDocType{MEMORANDO}
\SetDocTitle{Asunto del Memorando}
\SetDocCode{MEM-XXX-000}
\SetDocArea{Área Emisora}
\SetDocAuthor{Emisor}
\SetDocDate{\today}

\begin{document}

\MemoBlock%
  {Destinatario(s) — Cargo(s)}%
  {Nombre del emisor — Área}%
  {Asunto del memorando en una línea}%
  {MEM-XXX-000}

\section{Instrucción / Comunicación}

Texto principal del memorando. Usar lenguaje claro, directo y sin ambigüedades.
Un memorando estándar no debe superar una página.

\begin{enumerate}
  \item Primera instrucción o punto de comunicación.
  \item Segunda instrucción o punto de comunicación.
  \item Tercera instrucción o punto de comunicación, si aplica.
\end{enumerate}

\begin{ObservationBox}[Vigencia]
  Indicar desde cuándo aplica esta instrucción y si tiene fecha de vencimiento.
\end{ObservationBox}

\NoteInline{Acuse de recibo:} Confirmar recepción respondiendo a este documento
con copia al área emisora antes del \textbf{[fecha límite]}.

\SignatureBlock{Nombre Apellido}{Cargo del Emisor}

\ContactSignature

\end{document}
